\usepackage[letterpaper, margin=1in]{geometry}
\usepackage{enumitem}
\usepackage{mathtools}
\usepackage{amsthm}
\usepackage{amssymb}
\usepackage{mathrsfs}
\usepackage{xspace}
\usepackage{bm}
\usepackage[breakable]{tcolorbox}
\usepackage{longtable}
\usepackage{booktabs}
\usepackage{graphicx}
\usepackage{relsize}
% \usepackage[left]{lineno}
% \linenumbers
\usepackage{tikz}
  \usetikzlibrary{positioning,fit,calc}
  \tikzstyle{block} = [draw=black, ultra thin, text width=2cm, minimum height=1.3cm, font = {\footnotesize\itshape},align=center]  
  \tikzstyle{arrow} = [thick,->,>=stealth]
  
  \renewenvironment{quote}
  {\list{}{\leftmargin=0.5cm \rightmargin=0.5cm} % adjust margins here
   \item\relax}
  {\endlist}
  
    
\newcommand{\tree}{\mathrm{tree}}
\newcommand{\ttree}{\mathrm{ttree}}
\newcommand{\support}{\mathrm{support}}

  \usepackage[page,toc,titletoc,title]{appendix}
\usepackage[disable,textsize=footnotesize]{todonotes}
\usetikzlibrary{shapes,arrows,positioning,patterns}
\usepackage{comment}
\usepackage{hyperref}
\usepackage{cleveref}
\hypersetup{
    colorlinks=true,
    citecolor=blue,
    linkcolor=blue,
    filecolor=blue,      
    urlcolor=blue,
}
\usepackage{thm-restate}
\usepackage{placeins}
\usepackage[font=small,labelfont=bf,margin=0.06in]{caption}

%more flexibility with bib entries
\usepackage[style=alphabetic,maxnames=999]{biblatex}
\addbibresource{bib.bib}
\AtEveryBibitem{
  \clearfield{editor}
  \clearlist{editor}
  \clearname{editor}
  \clearfield{location}
  \clearlist{location}
  \clearname{location}
  \clearfield{isbn}
  \clearlist{isbn}
  \clearname{isbn}
  % \clearfield{url}
  % \clearlist{url}
  % \clearname{url}
  \ifentrytype{online}
    {}
    {\clearfield{url}}
}

\RequirePackage[osf,sc]{mathpazo}
% Margin kerning and slight font expansion, which avoid most overfull lines.
\usepackage{microtype}
% Extra stretch for the few paragraphs that still cannot be broken without
% an overfull line.  TeX uses it only in a final pass for such paragraphs,
% so all other paragraphs are unaffected.
\setlength{\emergencystretch}{1.5em}
%for vertical bars
\usepackage{framed}
\usepackage{mdframed}
\usepackage{float}
\usepackage{xcolor}



\newcommand{\ind}[2][]{%
  \mathrel{
    \mathop{
      \vcenter{
        \hbox{\oalign{\noalign{\kern-.3ex}\hfil$\vert$\hfil\cr
              \noalign{\kern-.7ex}
              $\smile$\cr\noalign{\kern-.3ex}}}
      }
    }^{#2}\displaylimits_{#1}
  }
}
\newcommand{\wh}{\widehat}
\newcommand{\nind}{\not\ind}
\renewcommand{\deg}{\mathrm{degeneracy}}
%note macros
\newcommand{\sz}[1]{\todo[color=pink!40]{sz: #1}}
\newcommand{\jan}[1]{\todo[color=green!40]{jan: #1}}
\newcommand{\niko}[1]{\todo[color=blue!20]{niko: #1}}
\newcommand{\nikoin}[1]{\todo[inline,size=\normalsize,color=blue!20,caption={Niko: [cannot render inline todos in listoftodos]}]{Niko: #1}}

\newcommand{\szfuture}[1]{\todo[color=green!40]{future: #1}}




\DeclareMathOperator{\rk}{rk}
\DeclareMathOperator{\VCdim}{VCdim}
\DeclareMathOperator{\TVCdim}{{2}VCdim}
\DeclareMathOperator{\fw}{fw}
\DeclareMathOperator{\cw}{cliquewidth}
\DeclareMathOperator{\rw}{rankwidth}
\DeclareMathOperator{\bfw}{bfw}
\DeclareMathOperator{\const}{const}
\newcommand{\fwp}{\fw^\mathrm{part}}%$k$-partite
\newcommand{\fwg}{\fw^\mathrm{gen}}%$k$-generated
\newcommand{\fwd}{\dfw}
% \newcommand{\fwd}{\fw^\mathrm{def}}%$k$-definable
\newcommand{\bfwp}{\bfw^\mathrm{part}}%$k$-partite
\newcommand{\bfwg}{\bfw^\mathrm{gen}}%$k$-generated
\newcommand{\bfwd}{\bfw^\mathrm{def}}%$k$-definable

\newcommand{\poly}{\mathrm{poly}}
\newcommand{\subpoly}{\mathrm{subpoly}}
\DeclareMathOperator{\iw}{iw}
\DeclareMathOperator{\copw}{copwidth}
\DeclareMathOperator{\dfw}{dfw}
\DeclareMathOperator{\adm}{adm}
\DeclareMathOperator{\imp}{imp}
\DeclareMathOperator{\tww}{tww}
\DeclareMathOperator{\wcol}{wcol}
\DeclareMathOperator{\scol}{scol}
\DeclareMathOperator{\sreach}{SReach}
\DeclareMathOperator{\Types}{Types}
\DeclareMathOperator{\wreach}{WReach}
\DeclareMathOperator{\loli}{loli}
\DeclareMathOperator{\cone}{Cone}
\DeclareMathOperator{\Ball}{Ball}

%%% LEFTBAR BEGIN %%%%%%%%%%%%%%%%%%%%%%%%%%%%%%%%%%%%%%%%%%%%%%%%%%%%%%%%%%%%%

%https://tex.stackexchange.com/questions/50612/alignment-in-the-leftbar-environment-is-rightbar-possible
\newlength{\leftbarwidth}
\setlength{\leftbarwidth}{3pt}
\newlength{\leftbarsep}
\setlength{\leftbarsep}{10pt}

\renewenvironment{leftbar}[1][blue]
{%
\def\FrameCommand
{%
{\hspace{-8pt} \color{#1} \vrule width 3pt}%
\hspace{0pt}%must no space.
\fboxsep=\FrameSep\colorbox{#1!5!}%
}%
\MakeFramed{\hsize\hsize\advance\hsize-\width\FrameRestore}%
}
{\endMakeFramed}
\setlength{\FrameSep}{5pt}
%%% LEFTBAR END %%%%%%%%%%%%%%%%%%%%%%%%%%%%%%%%%%%%%%%%%%%%%%%%%%%%%%%%%%%%%%%


\makeatletter
%define \szin environment, which which ignores its content if todonotes is used with "disabled" option, and shows a comment otherwise

  \if@todonotes@disabled % Check todnotes is used with disabled option
  \excludecomment{szin}
  \else %todonotes enabled
  \newenvironment{szin}{\begin{leftbar}[red]{\noindent\bf sz: }}{\end{leftbar}}
  \fi

  % Jan's inline notes, rendered like szin.
  \if@todonotes@disabled
  \excludecomment{janin}
  \else
  \newenvironment{janin}{\begin{leftbar}[blue!70!black]{\noindent\bf jan: }}{\end{leftbar}}
  \fi

  \if@todonotes@disabled % Check todnotes is used with disabled option
  \excludecomment{new}
  \else %todonotes enabled
  \newenvironment{new}{\begin{leftbar}[green]{\noindent\bf new: }}{\end{leftbar}}
  \fi

\makeatother


\newcommand\myitem[1]{\hyperref[item:#1]{\emph{\ref*{item:#1}.}{}}\xspace}

% Needed to give custom number to a theorem
% Usage: \begin{customthm}{\ref{thm:x}}[slightly reformulated]
\newtheorem{innercustomthm}{Theorem}
\newenvironment{customthm}[1]
  {\renewcommand\theinnercustomthm{#1}\innercustomthm}
  {\endinnercustomthm}

% \newtheorem{innercustomlem}{Theorem}
% \newenvironment{customlem}[1]
%   {\renewcommand\theinnercustomlem{#1}\innercustomlem}
%   {\endinnercustomlem}

% % Needed to give custom number to a corollary
% \newtheorem{innercustomcol}{Corollary}
% \newenvironment{customcol}[1]
%   {\renewcommand\theinnercustomcol{#1}\innercustomcol}
%   {\endinnercustomcol}

% Needed to give custom number to a definition
% \newtheorem{innercustomdef}{Definition}
% \newenvironment{customdef}[1]
%   {\renewcommand\theinnercustomdef{#1}\innercustomdef}
%   {\endinnercustomdef}
\newcommand{\appmark}{$\ast$}
%\theoremstyle{plain}
\newtheorem{theorem}{Theorem}[section]
% thm-restate (via aliasctr) overrides hyperref's \@addtoreset patch, so the
% hyperlink anchors would omit the section number and collide across sections.
% Set them by hand; this must precede the [theorem]-sharing environments,
% which copy \theHtheorem when they are declared.
\renewcommand{\theHtheorem}{\theHsection.\arabic{theorem}}
\newtheorem*{theorem*}{Theorem}
\newtheorem{theorem?}[theorem]{?`Theorem?}
\newtheorem{conjecture}[theorem]{Conjecture}
\newtheorem*{mcconjecture}{Model Checking Conjecture}
\newtheorem{question}[theorem]{Question}

\newtheorem{corollary}[theorem]{Corollary}
\newtheorem*{corollary*}{Corollary}
\newtheorem{lemma}[theorem]{Lemma}
\newtheorem*{lemma*}{Lemma}
\newtheorem{lemma?}[theorem]{?`Lemma?}
\newtheorem{fact}[theorem]{Fact}
\newtheorem{observation}[theorem]{Observation}
\newtheorem{notation}{Notation}
\newtheorem{proposition}[theorem]{Proposition}
\newtheorem*{proposition*}{Proposition}
\newtheorem{wrong}[theorem]{False Statement}
\newtheorem*{wrong*}{False Statement}

\renewcommand{\qedsymbol}{\raisebox{0.2ex}{\scalebox{0.85}{$\blacksquare$}}}
\crefname{theorem?}{Theorem{\scriptsize{?}}}{Theorems{\scriptsize{?}}}
\crefname{lemma?}{Lemma{\scriptsize{?}}}{Lemmas{\scriptsize{?}}}
\crefname{Theorem?}{Theorem{\scriptsize{?}}}{Theorems{\scriptsize{?}}}
\crefname{Lemma?}{Lemma{\scriptsize{?}}}{Lemmas{\scriptsize{?}}}


\newtheorem{claim}{Claim}[section]
\renewcommand{\theHclaim}{\theHsection.\arabic{claim}}
\crefname{claim}{claim}{Claims} % lowercase
\Crefname{claim}{Claim}{Claims} % capitalized

\newenvironment{claimproof}[1][\proofname]{%
  \renewcommand{\qedsymbol}{\raisebox{0.3ex}{\scalebox{1}{$\blacktriangleleft$}}}%
  \begin{proof}[#1]%
}{%
  \end{proof}%
}

\newtheorem{goal}[theorem]{Goal}
\newcommand{\msfw}{\mathrm{monotone\text{-}set\text{-}fw}}
\newcommand{\mfw}{\mathrm{monotone\text{-}fw}}


\theoremstyle{remark}
\newtheorem{remark}[theorem]{Remark}
\newtheorem{example}[theorem]{Example}
\newtheorem{definition}[theorem]{Definition}

% \numberwithin{equation}{section}

\def\cl{\mathrm{cl}}

\def\marker{\bar}
\def\bu{{\marker u}}
\def\bv{{\marker v}}
\def\bw{{\marker w}}
\def\bx{{\marker x}}
\def\by{{\marker y}}
\def\bz{{\marker z}}
\def\bs{{\marker s}}
\renewcommand{\L}{{\mathrm{L}}}
\newcommand{\G}{{\mathrm{G}}}

\newcommand{\parent}{\mathrm{parent}}
\renewcommand{\preceq}{\preccurlyeq}
\renewcommand{\succeq}{\succcurlyeq}


\def\Nesetril{Ne\v{s}et\v{r}il\xspace}
\def\Dvorak{Dvo\v{r}\'{a}k\xspace}
\def\Kral{Kr\'{a}l\xspace}
\def\Gajarsky{Gajarsk\'{y}\xspace}

\newcommand{\cmso}{\mathrm{CMSO}}
\newcommand{\mso}{\mathrm{MSO}}
\newcommand{\first}{\text{first}}
\newcommand{\str}[1]{{{#1}}}
\newcommand{\from}{\colon}
\newcommand{\tto}{\rightsquigarrow}
\newcommand{\pto}{\rightharpoonup}
\newcommand{\set}[1]{\{#1\}}
\newcommand{\mset}[1]{\{\!\!\{#1\}\!\!\}}
\newcommand{\setof}[2]{\set{#1 \,:\, #2 }}
\newcommand{\setoff}[2]{\set{#1 \,:\, #2 }}
\newcommand{\msetof}[2]{\mset{#1 \,|\, #2 }}
\def\phi{\varphi}
\def\cal{\mathcal}
\def\N{\mathbb N}
\def\R{\mathbb R}
\def\epsilon{\varepsilon}
\def\eps{\varepsilon}
\renewcommand{\subset}{\subseteq}
\renewcommand{\setminus}{-}
\renewcommand{\le}{\leqslant}
\renewcommand{\leq}{\le}
\renewcommand{\ge}{\geqslant}
\renewcommand{\geq}{\ge}
\newcommand{\lca}{\mathrm{lca}}
\newcommand{\dist}{\mathrm{dist}}
\newcommand{\lie}{\mathrm{lie}}
%treewidth and treedepth
\newcommand\tw{{\rm treewidth}}
\newcommand\td{{\rm treedepth}}
\renewcommand{\AA}{\cal A}
\newcommand{\BB}{\cal B}
\newcommand{\CC}{\mathscr C}
\newcommand{\DD}{\mathscr D}
\newcommand{\TT}{\cal T}
\newcommand{\HH}{\cal H}
\newcommand{\RR}{\cal R}
\renewcommand{\SS}{\cal S}
\newcommand{\FF}{\cal F}
\newcommand{\PP}{\cal P}
\newcommand{\UU}{\cal U}
\newcommand{\QQ}{\cal Q}
\newcommand{\trans}[1]{\mathsf{#1}}
% \newcommand{\merge}[2]{\stackrel{u\mapsto v}{\longrightarrow}}
\DeclareMathOperator{\mw}{mw}
\newcommand\reach{{\rm reach}}
\newcommand{\warp}[1]{\langle #1\rangle}
\newcommand{\leaves}{\mathrm{Leaves}}

\newcommand{\ATypes}{\mathrm{ATypes}}

% \def\H{\mathcal H}
\def\S{\mathcal S}
\newcommand{\Oof}{O}
\def\edge{\textit{edge}}
\newcommand{\tup}{\bar}
\newcommand{\tp}{\textnormal{tp}}
\newcommand{\ltp}{\textnormal{ltp}}
\newcommand{\ttp}{\textnormal{ttp}}
\newcommand{\atp}{\textnormal{atp}}
% \newcommand{\tpp}[1]{\textnormal{tp}_{#1}}
\def\typi{\textit{type}}
\newcommand{\stp}{\textnormal{stp}}
\newcommand{\stree}{\textnormal{stree}}
\newcommand{\ftree}{\textnormal{ftree}}
% \newcommand{\ltp}{\textnormal{ltp}}
% Part III (localglobal-trees.tex)
\newcommand{\distFO}{\textup{dist-FO}\xspace}
\newcommand{\merge}[2]{\mathrm{merge}_{#1,#2}}
\newcommand{\resolve}[3]{\mathrm{resolve}^{#1}_{#2,#3}}
\newcommand{\rhom}{\rho^-}
\newcommand{\rhop}{\rho^+}
\newcommand{\size}[1]{\lVert #1\rVert}
\newcommand{\crad}{R}  % clustering radius of lem:clustering and lem:to-ftree

% \newcommand{\dist}{\textnormal{dist}}
% \newcommand{\reach}{\textnormal{reach}}
\newcommand{\cluster}{\textnormal{cluster}}
\newcommand{\rep}{\textnormal{rep}}
\newcommand{\FlipBudget}{\textsf{Flips}}
\newcommand{\sample}{\mathrm{sample}}


\newcommand{\anonym}[2][]{#2}%for anonymized version put #1, for non-anonymized put #2

%%%%%%%%%%%%%%%%%%%%%%%%%%%%%%%%%%%%%%%%%%%%%%%%%%%%%%%%%%%%%%%%%%%%%%%%%%%%%%%
%%%%%%%%%%%%%%%%%%%%%%%%%%%%%%%%%%%%%%%%%%%%%%%%%%%%%%%%%%%%%%%%%%%%%%%%%%%%%%%

% \newcommand{\ERCagreement}{This paper is part of a project that has received funding from the European Research Council (ERC) (grant agreement No 101126229 -- {\sc BUKA}).}


\newcommand{\ERCagreement}{\xspace ST received funding from the European Research Council (ERC) (grant agreement №948057 -- {\sc bobr} -- and №101126229 -- {\sc buka}).
% \begin{picture}(0,0)
%   \put(412,-430)
%   {\hbox{\includegraphics[width=40px]{logo-erc.jpg}}}
%   \put(402,-490)
%   {\hbox{\includegraphics[width=60px]{logo-eu.pdf}}}
%   \end{picture}%
\begin{tikzpicture}[remember picture, overlay]
  % Place the ERC logo in the lower-right corner
  \coordinate (refpoint) at (12,0); %(current page.south east); 

  \node[anchor=south, yshift=0cm] at (refpoint)
  {\includegraphics[width=40px]{logo-erc.jpg}};

  % Place the EU logo slightly below the ERC logo
  \node[anchor=south, yshift=-2cm] at (refpoint)
  {\includegraphics[width=60px]{logo-eu.pdf}};
\end{tikzpicture}
}
